%!TEX root = ../template.tex
%%%%%%%%%%%%%%%%%%%%%%%%%%%%%%%%%%%%%%%%%%%%%%%%%%%%%%%%%%%%%%%%%%%%
%% glossary-xltabular.tex
%% NOVA thesis documentation file
%%
%% Detailed explanation of glossary-xltabular.sty functionality.
%%%%%%%%%%%%%%%%%%%%%%%%%%%%%%%%%%%%%%%%%%%%%%%%%%%%%%%%%%%%%%%%%%%%

\typeout{NT FILE glossary-xltabular.tex}

\chapter{Explanation of glossary-xltabular.sty}
\label{annex:glossary-xltabular}

The file \texttt{NOVAthesisFiles/StyFiles/glossary-xltabular.sty} is a custom extension for the \texttt{glossaries} package. It provides new glossary styles that use the \texttt{xltabular} package, which combines the features of \texttt{tabularx} (flexible-width columns) and \texttt{longtable} (multi-page tables).

\section{Glossary Styles}
The package defines several sophisticated styles for presenting glossaries, acronyms, and symbols in a professional, tabular format:

\begin{itemize}
    \item \textbf{xltabular}: A basic two-column layout where the description column (the \texttt{X} column) expands to fill the available page width.
    \item \textbf{xltabular3col}: A three-column layout (name, description, page list) where the description is automatically sized.
    \item \textbf{xltabular4col}: A four-column layout, typically used for symbols, which includes name, description, symbol, and page list.
    \item \textbf{altxltabular4col}: A variation of the four-column style that adjusts which columns are fixed-width and which are flexible.
\end{itemize}

\section{Key Features}
Each of the styles listed above also comes with variants that add specific formatting:
\begin{itemize}
    \item \textbf{header}: Adds a bold header row (e.g., ``Name'', ``Description'') to the top of the table on every page.
    \item \textbf{border}: Adds horizontal and vertical borders (lines) to create a formal grid.
\end{itemize}

\section{Layout Customization}
A unique feature in the \emph{NOVAthesis} version of this style is its integration with the template's configuration system. The \texttt{theglossary} environment dynamically pulls its column layout from the template's settings (\texttt{/novathesis/glossaries/layout}), allowing users to change the look of their glossary simply by adjusting a key-value option in the config files.

\section{Summary}
In summary, \texttt{glossary-xltabular.sty} provides high-quality, multi-page tabular layouts for lists of acronyms and symbols, ensuring they are both visually appealing and properly aligned across the entire document.
